\section{Conditional robustness and discrepancy learning}\label{sec:C}
\subsection{Posterior component odds}\label{sec:C.1}
Fix a discrepancy leaf, the other tree contributions, $n\geq1$, $\sigma_1^2>0$, $0<w<1$ and ordered positive variances $\tau_0^2<\tau_1^2$. Let $\bar r$ be its trial-control partial-residual mean. Bayes' rule applied to the two normal marginal densities gives the posterior spike odds
\begin{equation}\label{eq:odds}
\frac{P(z=1\mid\mathbf r)}{P(z=0\mid\mathbf r)}
=\frac{w}{1-w}\sqrt{\frac{\tau_1^2+\sigma_1^2/n}{\tau_0^2+\sigma_1^2/n}}
\exp\left[-\frac{\bar r^2}{2}
\left\{\frac{1}{\tau_0^2+\sigma_1^2/n}-\frac{1}{\tau_1^2+\sigma_1^2/n}\right\}\right].
\end{equation}
For this proof only, denote the positive leading coefficient by $C$ and the positive coefficient of $\bar r^2$ in the exponent by $\kappa$. Then the spike probability is $\gamma(\bar r)\coloneqq P(z=1\mid\mathbf r)=C e^{-\kappa\bar r^2}/(1+C e^{-\kappa\bar r^2})$. 

\subsection{Proof of main-text Theorem 2}\label{sec:C.2}
The posterior mean under component $k$ is $\omega_k\bar r$, where $\omega_k\coloneqq \tau_k^2/(\tau_k^2+\sigma_1^2/n)$. Because $\omega_1>\omega_0$, the difference between the mixture posterior mean and the slab-only posterior mean is
\begin{equation}
\Delta_{\mathrm{mix}}(\bar r)\coloneqq -\gamma(\bar r)(\omega_1-\omega_0)\bar r.
\end{equation}
Using $\gamma(\bar r)\leq C e^{-\kappa\bar r^2}$,
\begin{equation}
|\Delta_{\mathrm{mix}}(\bar r)|\leq(\omega_1-\omega_0)C|\bar r|e^{-\kappa\bar r^2}
\leq\frac{(\omega_1-\omega_0)C}{\sqrt{2e\kappa}}.
\end{equation}
The first upper bound tends to zero in both tails, proving boundedness and tail attenuation. This is the extra shrinkage attributable to the spike relative to the slab-only mean. Relative to the unshrunk likelihood mean $\bar r$, the remaining slab shrinkage is $(\omega_1-1)\bar r$, which is unbounded for fixed $n$ and finite $\tau_1^2$. The result bounds additional mixture shrinkage relative to the slab, not total shrinkage relative to the likelihood mean.

\subsection{Proof of main-text Theorem 3}\label{sec:C.3}
Equation \eqref{eq:odds} is strictly decreasing in $|\bar r|>0$. If $C>1$, the posterior spike probability exceeds one half exactly when $|\bar r|<\sqrt{\log(C)/\kappa}$; if $C\leq1$, it never exceeds one half. There is no universal threshold independent of $w$, sample size and scales.

Now suppose the partial residuals are independent $\N(d,\sigma_1^2)$ observations. Almost surely $\bar r\to d$. For both $k=0,1$, the conditional posterior mean $\omega_k\bar r\to d$ and variance $(n/\sigma_1^2+\tau_k^{-2})^{-1}\to0$. Each normal component therefore concentrates at $d$, and so does every convex mixture of them. If $|d|\ne c$, its posterior probability of $(-c,c)$ converges to $\ind\{|d|<c\}$.

Nevertheless, the posterior spike probability converges to
\begin{equation}
\frac{w\,\phi(d;0,\tau_0^2)}{w\,\phi(d;0,\tau_0^2)+(1-w)\,\phi(d;0,\tau_1^2)},
\end{equation}
which lies strictly between zero and one for finite $d$ under the stated assumptions. The overlapping component densities therefore prevent consistent identification of a unique mixture allocation at a finite leaf mean. The boundary $|d|=c$ is excluded from the probability-limit claim.

\subsection{Scope of the conditional assertions}\label{sec:C.4}
The proofs fix the tree partition and the other tree contributions. In the full model, those contributions, the partition, $w$ and the spike variance are estimated and dependent. Posterior concentration of their sum requires additional assumptions and arguments. A sum of tree functions also admits multiple tree decompositions, so identifiability of individual terminal-node effects is not the same as identification of $f$ or $f+g$.

With both control sources and adequate support, the historical likelihood informs $f$ and the trial-control likelihood informs $f+g$. The resulting difference can be learned where both functions are estimable. Without trial controls, the likelihood contains no information about $g$ under the separate treatment model. Neither the leaf theorem nor an ESS calculation supplies that missing information. The regional-discrepancy figure reports finite-sample empirical regional behavior; Table 7 reports prior information and is not evidence for conditional-discrepancy consistency in the single-arm application.
